% 可読性のための工夫の見本．
% 既定のままでも泣き別れ抑制などは働くが，紙面の見た目を変える機能は
% すべてオプトインである．ここでは柱・段間罫・和欧間アキを有効にする．
\documentclass[lualatex,paper=b5j,fontsize=10pt,twoside]{jlreq}

\usepackage[
  runninghead,
  columnrule,
  interspacing
]{surikumi}

\MathMagazineSetup{
  feature={},
  series={要点の整理／数学 B},
  series-position=left,
  pattern=origami,
  title={読みやすさの工夫},
  author={数学編集部},
  author-font=mincho,
  author-position=right,
  title-fill=black,
  deck={高密度の紙面を保ったまま読みやすくするための部品をまとめる．}
}

\begin{document}

\MakeMathMagazineTitle
\MMConfirmationReset

\MMConfirmationSection{強調のしかた}

\MMLead{本文と同じ書体・同じ大きさのまま，読者の目を止める方法を三つ用意する．
太字や色に頼らず，和文組版の作法に沿った手段を選ぶ．}

\MMConfirmationSubsection{圏点による強調}

文章の中で一時的に声を強めたいところには\MMKenten{圏点}を打つ．
圏点は本文の書体・ウェイト・大きさを変えないので，
$8$ ポイントの高密度な紙面でも行の色が乱れない．
たとえば「和の上端と下端は\MMKenten{必ず}確認する」のように使う．

\MMConfirmationSubsection{用語の初出}

\MMTerm{階差数列}のように，用語の初出だけをゴシックで示す．
以後は本文の明朝体のままでよい．
\MMTerm{不動点}，\MMTerm{母関数}なども同じ扱いにする．

\MMConfirmationSubsection{ルビ}

読みにくい術語には\MMRuby{漸化式}{ぜんかしき}のようにルビを添える．
語ごとに区切るときは\MMRuby{階|差|数|列}{かい|さ|すう|れつ}と縦棒で対応させる．

\MMConfirmationSection{公式とまとめ}

\MMLead{公式の提示と答の強調は別物である．
答・最終結果を四角い枠で囲んではならないという規約は変わらない．}

\MMFormulaBox{%
  \[
    \sum_{k=1}^{n}k=\dfrac{n(n+1)}{2},\qquad
    \sum_{k=1}^{n}k^2=\dfrac{n(n+1)(2n+1)}{6}
  \]}

上の囲みは閉じた四角ではなく，左の太い罫と淡い地色でできている．
そのため，答を囲んだ枠と見間違えることがない．

\MMConfirmationGeneralComment{和の公式は，つねに上端と下端を先に読む．
$\sum_{k=1}^{n}$ と $\sum_{k=0}^{n-1}$ は別物である．}

\MMSummary{等差数列も等比数列も，一般項と和は初項と公差（公比）だけで決まる．
和の公式は上端と下端の確認から始める．}

\MMConfirmationSection{図表と全段幅}

\MMConfirmationSubsection{キャプション}

図と表には別々の番号が付く．

\begin{center}
  \begin{tikzpicture}[x=1mm,y=1mm]
    \draw[draw=SMrule,line width=.45pt] (0,0)--(46,0);
    \foreach \i in {0,1,2,3,4}{%
      \fill[SMink] ({4+9*\i},0) circle (.7);
      \node[below=1.2mm,font=\fontsize{6}{7}\selectfont]
        at ({4+9*\i},0) {$a_{\the\numexpr\i+1\relax}$};
    }
  \end{tikzpicture}
\end{center}
\MMCaption{等差数列の項は等間隔に並ぶ}

表のキャプションは表の上に置く．

\MMCaption[table]{よく使う和の公式}
\begin{center}
  \fontsize{7}{9}\selectfont
  \begin{tabular}{@{}ll@{}}
    \hline
    数列 & 和 \\
    \hline
    等差 & $S_n=\dfrac{a_1+a_n}{2}n$ \\[.6em]
    等比 & $S_n=\dfrac{a_1(1-r^n)}{1-r}$ \\
    \hline
  \end{tabular}
\end{center}

\MMConfirmationSubsection{一時的に全段幅を使う}

$1$ 段に収まらない式は，式を分けるか，次のように一時的に全段幅へ逃がす．

\begin{mmwideblock}
  \[
    \sum_{k=1}^{n}k(k+1)(k+2)(k+3)
      =\dfrac{n(n+1)(n+2)(n+3)(n+4)}{5},\qquad
    \sum_{k=1}^{n}k^4=\dfrac{n(n+1)(2n+1)(3n^2+3n-1)}{30}
  \]
\end{mmwideblock}

\MMConfirmationSection{行分割と改ページ}

\MMLead{既定で有効な工夫は目に見えない．段の頭や末尾に取り残される
$1$ 行をなくし，狭い段でも行があふれないようにする．}

段の末尾に段落の第 $1$ 行だけが残ること，段の先頭に段落の最終行だけが
置かれることを，いずれも禁じている．長い \verb|align*| は段をまたげるので，
段の下に大きな空きが生まれにくい．行内の長い式は関係記号のあとで折り返せる．

たとえば $\left(\sum_{k=1}^{n}a_kb_k\right)^2\leq
\left(\sum_{k=1}^{n}a_k^2\right)\left(\sum_{k=1}^{n}b_k^2\right)$
のような式は，狭い段でも $\leq$ の位置で分割できる．

\MMConfirmationGeneralNote{これらは \texttt{penalties} オプションで
まとめて無効にできる．}

\MMConfirmationSection{読者が疲れにくい版}

\MMLead{長時間読ませる記事には，第三の密度 \texttt{density=airy} を使う．
行送り，別行数式の前後，\texttt{align} の行間がまとめて開く．}

\begin{verbatim}
\usepackage[density=airy,interspacing]{surikumi}
\end{verbatim}

既存の \texttt{compact} と \texttt{standard} の値は変わらないので，
既刊の紙面は影響を受けない．

\MMConfirmationSection*{付録：柱と段間罫の確認}

以下は，柱と段間罫の見え方を確かめるための本文である．
柱は記事の扉には出ず，$2$ ページ目以降に出る．
段間罫は段の高さに合わせて引かれるため，段が短いページでは短くなる．

数列 $\{a_n\}$ が\MMTerm{等差数列}であるとは，$a_{n+1}-a_n$ が $n$ に
よらず一定であることをいう．この一定値を公差といい，$d$ で表す．
一般項は $a_n=a_1+(n-1)d$ であり，和 $S_n=\sum_{k=1}^{n}a_k$ は
初項と末項の平均に項数を掛けて
\[
  S_n=\dfrac{a_1+a_n}{2}n
\]
となる．末項 $a_n$ を代入すれば $S_n=\dfrac{2a_1+(n-1)d}{2}n$ とも書ける．

\MMTerm{等比数列}では $a_{n+1}/a_n$ が一定であり，この値を公比という．
一般項は $a_n=a_1r^{n-1}$ である．和は $r\neq1$ のとき
\[
  S_n=\dfrac{a_1(1-r^n)}{1-r}
\]
であり，$r=1$ のときは $S_n=na_1$ となる．
$r=1$ を分母に代入できないので，\MMKenten{必ず}場合を分ける．

階差数列を $b_n=a_{n+1}-a_n$ とおくと，$n\geq2$ のとき
\[
  a_n=a_1+\sum_{k=1}^{n-1}b_k
\]
が成り立つ．この式は $n=1$ では使えないため，$n=1$ の場合を別に確かめる．
和から一般項を求めるときも同じで，$a_n=S_n-S_{n-1}$ は $n\geq2$ でしか
使えない．$n=1$ では $a_1=S_1$ を用いる．

\MMSummary{一般項の公式も和の公式も，使える $n$ の範囲がある．
範囲の端では，必ずもとの定義に戻って確かめる．}

\end{document}
